% !TeX root = linear-codes.tex
\section{Introduction}
\label{sec:intro}
In this paper we exhibit binary linear codes of blocklength
$n$ and size $\Omega(n)$ times the Gilbert--Varshamov
bound for any distance $d$ in the range
$\lg n\ll d\le(1/2-\eps)n$. Denote
\begin{align*}
  \ball=\setbuilder{x\in\field_2^n}{0<|x|<d}
\end{align*}
the set of all 0-1 vectors with weight between 1 and $d-1$,
and set $V=|\ball|$. The Gilbert--Varshamov bound
\cite{Gilbert1952, Varshamov1957} says that there exists a
code of size at least
\begin{align*}
  \frac{2^n}{V+1},
\end{align*}
and such a code may be chosen to be linear
\cite{Varshamov1957}. For fixed relative distance $d/n$,
Jiang and Vardy \cite{JiangV2004} obtained a superconstant
multiplicative improvement by exhibiting nonlinear codes
whose size is $\Omega(\lg V)$
times the Gilbert--Varshamov bound. The improvement is
$\Omega(n)$ when $d/n$ is a constant and $\Omega(d\lg(n/d))$
for $d\ll n$.

Gaborit and Zémor \cite{GaboritZ2008} obtained an
$\Omega(n)$ improvement for linear codes, but only at rate
$1/2$ and for certain blocklengths $n=2p^m$. The rate
parameter $1/2$ corresponds to distance
$d=\delta_0n+\Theta(\lg n)$, where $\delta_0\approx0.11$ is
the constant satisfying $\entb{\delta_0}=1/2$. More
recently, Hao, Huang, Livshyts, and Tikhomirov
\cite{HaoHLT2022} obtained an $\Omega(\sqrt n)$ improvement
for general binary linear codes at fixed relative distance
$d/n$. Yuan and Zhu \cite{YuanZ2026} gave another proof of
the same $\Omega(\sqrt n)$ improvement and extended the
result to quantum codes.
%
% !TeX root = linear-codes.tex
% The binary matrix-vector product illustrated with marks and empty cells.
\begin{center}
\begin{tikzpicture}[x=2.3mm,y=2.3mm]
  \tikzset{binary grid line/.style={
    draw=black!55,
    line width=.22pt,
    line cap=butt
  }}
  \def\zeromarkhalfwidth{.36}
  \def\zeromarkcorner{.90pt}

  % Draw a binary grid from row-major data, with the top row first.
  % Arguments: lower-left x, lower-left y, rows, columns, data.
  % Ones are centered purple rounded squares that do not touch the cell walls.
  % The complete grid is drawn afterward as
  % rows+1 uninterrupted horizontal lines and columns+1 vertical lines.
  \newcommand{\drawbinarygrid}[5]{%
    \path[fill=white] (#1,#2) rectangle ++(#4,#3);
    \edef\binarydata{#5}
    \foreach \bit [count=\cell from 0] in \binarydata {%
      \pgfmathtruncatemacro{\gridrow}{int(\cell/#4)}
      \pgfmathtruncatemacro{\gridcol}{mod(\cell,#4)}
      \ifnum\bit=1\relax
        \path[fill=purple,rounded corners=\zeromarkcorner]
          ({#1+\gridcol+.5-\zeromarkhalfwidth},
           {#2+#3-\gridrow-.5-\zeromarkhalfwidth})
          rectangle
          ({#1+\gridcol+.5+\zeromarkhalfwidth},
           {#2+#3-\gridrow-.5+\zeromarkhalfwidth});
      \fi
    }
    \foreach \gridrow in {0,...,#3} {
      \draw[binary grid line]
        (#1,{#2+\gridrow}) -- ++(#4,0);
    }
    \foreach \gridcol in {0,...,#4} {
      \draw[binary grid line]
        ({#1+\gridcol},#2) -- ++(0,#3);
    }
  }

  \def\matrixrows{10}
  \def\matrixcols{18}
  \def\matrixleft{7}
  \def\matrixbottom{1}

  \def\xdata{0,1,0,0,1,0,0,0,0,0,0,0,0,1,0,1,0,0}
  \def\ydata{0,0,0,0,0,0,0,1,1,1}
  \def\matrixdata{
    1,1,0,1,0,0,1,1,0,0,1,1,0,1,1,0,0,1,
    0,0,1,0,1,1,1,0,1,1,1,0,1,0,0,1,1,0,
    0,0,1,1,0,1,0,1,1,0,0,1,1,0,1,0,1,1,
    1,1,0,0,1,0,1,1,0,1,1,0,0,1,0,1,0,0,
    1,1,1,0,0,1,0,0,1,1,0,1,1,0,0,1,1,0,
    0,0,1,1,1,0,1,1,1,0,1,0,0,1,1,0,0,1,
    1,1,1,0,0,0,0,1,0,1,1,0,1,0,0,1,1,0,
    0,0,0,1,1,1,1,0,1,0,1,1,0,0,1,0,0,1,
    0,1,0,1,1,0,1,0,1,1,0,1,0,1,1,0,1,1,
    1,1,1,0,0,0,1,1,0,1,0,0,1,1,0,1,0,0
  }

  % Derive the layout from the matrix dimensions.
  \pgfmathsetmacro{\matrixcenterx}{\matrixleft+\matrixcols/2}
  \pgfmathsetmacro{\matrixcentery}{\matrixbottom+\matrixrows/2}
  \pgfmathsetmacro{\matrixtop}{\matrixbottom+\matrixrows}
  \pgfmathsetmacro{\xgridbottom}{\matrixtop+1.5}
  \pgfmathsetmacro{\xgridtop}{\xgridbottom+1}
  \pgfmathsetmacro{\legendheight}{\matrixbottom-.45}

  % x is aligned with the matrix columns and separated by 1.5 cells.
  \drawbinarygrid{\matrixleft}{\xgridbottom}
    {1}{\matrixcols}{\xdata}
  \drawbinarygrid{0}{\matrixbottom}
    {\matrixrows}{1}{\ydata}
  \drawbinarygrid{\matrixleft}{\matrixbottom}
    {\matrixrows}{\matrixcols}{\matrixdata}

  \node[left=2.5pt] at (0,\matrixcentery) {$y$};
  \node at (3.8,\matrixcentery) {$=$};
  \node[above=1.5pt] at (\matrixcenterx,\xgridtop) {$x$};
\end{tikzpicture}

\vspace{2mm}
{\footnotesize The filled cells of $x$ select columns of $M$;
their sum is $y=Mx$.}
\end{center}

%

We show that a straightforward refinement of the folklore
union-bound argument gives an $\Omega(n)$ multiplicative
improvement to the Gilbert--Varshamov bound uniformly
throughout the range $\lg n\ll d\le(1/2-\eps)n$. This
establishes the largest known linear codes for
$e \sqrt{n} + \Omega(\lg n)\le d \le (1/2-\eps)n$.
%
Recall that a linear code of size $2^{n-r}$ may be specified
through its parity check matrix $H\in\field_2^{r\times n}$
satisfying $Hx \ne 0$ for all $x\in \ball$. Let $M$ be
chosen uniformly at random from $\field_2^{r\times n}$ and
put $\kero=\setbuilder{x\in \ball}{Mx = 0}$. By linearity of
expectation,
%
$\E\sparen{\size{\kero}}=\size{\ball}2^{-r}$.
%
When $r>\lg\size{\ball}$, we get
$\Pr\sparen{\kero\neq\emptyset} \le \E\sparen{\size{\kero}}<
1$, therefore an $r\times n$ parity check matrix for a
distance $d$ code is guaranteed to exist, which is the
folklore union-bound argument.

If we choose $r = \lg\size{\ball}-\lg(cn)$ instead, we have
$\E[\size{\kero}]=cn$ and we can no longer use the union
bound to argue $\Pr[\kero=\emptyset]>0$. We show that, for
small enough $c$, there is a choice of $M$ for which $\kero$
contains no odd circuit, that is, no inclusion-minimal
linearly dependent subset of odd cardinality. This allows
us to extend $M$ with just one more carefully chosen
parity-check row which maps all elements of $\kero$ to
nonzero. Indeed, every dependency among the elements of
$\kero$ then has even cardinality. Thus, if $\mathbf K$ is
the matrix whose columns are the elements of $\kero$, the
all-ones row lies in its row space, and
$v\mathbf K=\mathbf{1}^\transpose$ has a solution.
Extending $M$ with such a row $v$ gives the desired matrix.
A slight modification also proves linear independence of
$\kero$, with slightly worse intermediate constants; see
Remark~\ref{rem:linear-independence}.
